\chapter{Deux frères}
\label{chap:deux-freres}

Les mois avaient passé et la présence des deux garçons au sein de la caravane
n'avait plus rien d'exceptionnel. Ils avaient grandi avec elle et approchaient
désormais de leurs quatorze ans. Ils participaient à la vie du convoi comme les
autres, connaissaient les habitudes de chacun et avaient eux-mêmes développé
les leurs. On les voyait souvent ensemble, parfois occupés à travailler,
parfois à discuter, parfois à se chamailler pour des raisons qui semblaient
suffisamment importantes pour eux et parfaitement insignifiantes pour tous les
autres.

Un soir, alors que le camp venait d'être installé, un nouveau voyageur se
retrouva à discuter avec les deux garçons. La conversation n'avait rien de
mémorable : quelques questions sur la route, la caravane et leur prochaine
destination. Puis, comme il ne les connaissait pas encore, l'homme finit
naturellement par demander leurs noms. Kamatsu se présenta le premier.

Lorsque vint son tour, Khaer Ghal répondit sans réfléchir :

— Kael.

Kamatsu tourna légèrement la tête vers lui. Ce n'était qu'un mouvement presque
imperceptible, accompagné d'un petit sourire que leur interlocuteur ne pouvait
pas comprendre. Khaer Ghal croisa son regard et comprit immédiatement pourquoi.

Kael.

Il n'avait pas hésité. C'était la première fois que, lorsqu'on lui demandait
son nom, il répondait ainsi plutôt que de donner le nom orc qu'il avait porté
jusque-là. Il ne l'avait pas décidé et n'avait pas réfléchi à la réponse qu'il
devait donner. Le nom était simplement venu, naturellement, comme s'il avait
depuis longtemps cessé d'être un simple surnom donné par Kamatsu pour devenir
le sien.

Le sourire de son ami lui fit prendre pleinement conscience de ce qui venait
de se passer. Ils échangèrent un regard et, pendant un instant, semblèrent
partager exactement la même pensée. Aucun ne fit pourtant de commentaire. Il
n'y en avait pas besoin.

\illustration
  {1300.png}
  {Deux frères}
  {deux-freres}

L'homme les observa, amusé par cet échange silencieux.

— À vous voir comme ça, on pourrait vous croire frères.

La remarque était anodine, presque lancée en passant. Pourtant, elle sembla
suspendre un instant la conversation. Les deux garçons se regardèrent à
nouveau. Cette fois, il n'y eut ni plaisanterie ni réponse immédiate, seulement
ce même regard complice, un peu plus long que le précédent.

Ils ne partageaient ni le même sang, ni les mêmes parents, ni même les mêmes
origines. Quelques années auparavant, rien n'aurait dû permettre à leurs
chemins de se croiser : l'un était né parmi les Orcs, l'autre parmi les
humains, et leur première rencontre avait eu lieu de part et d'autre d'une
barrière, lorsque l'un était encore membre du peuple qui retenait l'autre
prisonnier.

Depuis, leurs anciennes vies avaient été bouleversées. Ils avaient survécu
ensemble et fui ensemble. Lorsque les caravaniers avaient cru Kamatsu
prisonnier de Khaer Ghal, il avait été le premier à prendre sa défense. Ils
avaient ensuite partagé les routes, les repas, les corvées, les découvertes et
les disputes.

Ils s'étaient inquiétés ensemble pour Bartiméus, avaient appris côte à côte
auprès de ceux que la caravane plaçait sur leur chemin et continuaient, avec un
sérieux très relatif, à s'entraîner l'un contre l'autre dès qu'ils en avaient
l'occasion. Sans jamais avoir besoin d'en parler, chacun était devenu pour
l'autre la personne qui restait lorsque tant d'autres étaient partis.

Le voyageur, qui ignorait évidemment tout de ce qui venait de traverser leur
esprit, poursuivit simplement la conversation. Kael et Kamatsu échangèrent
encore un sourire avant de reporter leur attention sur lui. Ce n'était pas une
décision et ils ne devenaient rien de nouveau à cet instant. Ils venaient
simplement de reconnaître ce qu'ils étaient devenus depuis longtemps sans
jamais avoir pensé à lui donner un nom.

La conversation continua sur autre chose et les deux garçons n'en reparlèrent
pas.

Ils n'en avaient pas besoin.

À partir de ce jour, rien ne changea réellement entre eux. Ils continuèrent à
se chamailler, à apprendre, à partir chacun de leur côté avant de se retrouver
autour du même feu. Kamatsu resta Kamatsu et Kael resta Kael.

Mais quelque chose avait malgré tout trouvé son nom.

Ils n'étaient plus seulement deux enfants que les circonstances avaient réunis,
ni deux amis qui voyageaient ensemble. Ils étaient devenus frères, non par le
sang, mais par tout ce qu'ils avaient traversé et tout ce qu'ils avaient choisi
de continuer à partager.

Des frères de cœur.